\documentclass[10pt]{beamer}
\usetheme{metropolis}
\usepackage{graphicx}
\usepackage{listings}
\usepackage{booktabs}
\usepackage{xcolor}
\definecolor{codebg}{HTML}{F5F5F5}
\definecolor{codegreen}{HTML}{2E7D32}
\definecolor{codegray}{HTML}{757575}
\definecolor{codeblue}{HTML}{1565C0}
\lstdefinestyle{rstyle}{language=R,backgroundcolor=\color{codebg},basicstyle=\ttfamily\scriptsize,
  keywordstyle=\color{codeblue}\bfseries,stringstyle=\color{codegreen},
  commentstyle=\color{codegray}\itshape,breaklines=true,frame=single,
  rulecolor=\color{codegray!30},showstringspaces=false,tabsize=2,
  morekeywords={library,ggplot,aes,geom_point,scale_color_viridis,scale_size,
    scale_shape_manual,theme_minimal,labs,ggMarginal,ggExtra}}
\lstdefinestyle{pythonstyle}{language=Python,backgroundcolor=\color{codebg},basicstyle=\ttfamily\scriptsize,
  keywordstyle=\color{codeblue}\bfseries,stringstyle=\color{codegreen},
  commentstyle=\color{codegray}\itshape,breaklines=true,frame=single,
  rulecolor=\color{codegray!30},showstringspaces=false,tabsize=4,
  morekeywords={import,from,as,pd,np,plt,sns,scatter,jointplot,hexbin}}
\title{High-Dimensional Data: Challenges \& Strategies}
\subtitle{Workshop 5 --- Module 1}
\author{Prof.Asoc.\ Endri Raco}
\institute{Data Visualization for Data Scientists \\ UNYT Tirana}
\date{2025/26}
\begin{document}

\begin{frame}
  \titlepage
\end{frame}

\begin{frame}{Module Roadmap}
  \begin{enumerate}
    \item The curse of dimensionality: why high-D data is hard to plot
    \item Visual encoding channels ranked by perceptual accuracy
    \item Multi-encoded scatter: 5 variables on one chart
    \item Aesthetic overload: when encoding fails
    \item Five strategies for high-dimensional visualization
    \item Aggregation techniques for large $n$
  \end{enumerate}
  \begin{exampleblock}{Learning Outcome}
    You will understand why standard 2D plots fail for high-dimensional
    data, apply multiple encoding channels without overload, and choose
    the right strategy (reduce, encode, facet, parallel, cluster) for
    any multivariate analysis.
  \end{exampleblock}
\end{frame}

\section{The Curse of Dimensionality}

\begin{frame}{Same $n$, More Dimensions $\rightarrow$ Sparser Data}
  \begin{center}
    \includegraphics[width=0.95\textwidth]{figures_w05m01/curse}
  \end{center}
  \begin{alertblock}{Pro-Tip}
    With $p$ dimensions and $n$ points, the proportion of the
    unit hypercube covered falls as $n^{-1/p}$. With 20 points in
    10D, each point is essentially alone --- there are no ``neighbours''
    to compare it with. This is why we need \textbf{dimension reduction}
    or \textbf{clever encoding} to visualize high-D data.
  \end{alertblock}
\end{frame}

\begin{frame}{What Changes at High Dimensions?}
  \centering
  \begin{tabular}{lp{4cm}p{4cm}}
    \toprule
    \textbf{Dimension} & \textbf{Standard Approach} & \textbf{High-D Challenge} \\
    \midrule
    1 variable & Histogram, boxplot & Still works (marginally) \\
    2 variables & Scatter plot & Still works \\
    3 variables & Scatter + colour/size & Starts getting cluttered \\
    5--10 variables & Pairs plot ($p^2$ panels) & Becomes unwieldy \\
    10--100 variables & ??? & Need dimension reduction \\
    100+ variables & ??? & Need algorithmic reduction + summary vis \\
    \bottomrule
  \end{tabular}
\end{frame}

\section{Visual Encoding Channels}

\begin{frame}{Ranked by Perceptual Accuracy}
  \begin{center}
    \includegraphics[width=0.92\textwidth]{figures_w05m01/encoding_channels}
  \end{center}
\end{frame}

\begin{frame}{Encoding Budget: How Many Variables Per Chart?}
  \centering
  \begin{tabular}{lp{5cm}r}
    \toprule
    \textbf{Encoding} & \textbf{Channel} & \textbf{Effective Capacity} \\
    \midrule
    Position x & Horizontal axis & 1 continuous \\
    Position y & Vertical axis & 1 continuous \\
    Colour hue & Categorical groups & 1 categorical ($\leq$8 levels) \\
    Colour value & Continuous intensity & 1 continuous (limited precision) \\
    Size & Bubble diameter & 1 continuous (3--5 distinguishable sizes) \\
    Shape & Point marker & 1 categorical ($\leq$6 shapes) \\
    Facet & Small multiples & 1--2 categorical \\
    \midrule
    \textbf{Total per chart} & & \textbf{5--7 variables max} \\
    \bottomrule
  \end{tabular}
  \vspace{6pt}
  {\small Beyond 5 variables on a single chart, add \textbf{faceting}
  or switch to a fundamentally different chart type (parallel coordinates,
  heatmap, PCA).}
\end{frame}

\section{Multi-Encoded Scatter}

\begin{frame}{Five Variables on One Scatter Plot}
  \begin{center}
    \includegraphics[width=0.72\textwidth]{figures_w05m01/multi_scatter}
  \end{center}
\end{frame}

\begin{frame}[fragile]{Building the Multi-Encoded Scatter}
  \begin{columns}[T]
    \begin{column}{0.48\textwidth}
      \textbf{R}
\begin{lstlisting}[style=rstyle]
ggplot(ecar, aes(
    x = FICO,           # position x
    y = Rate,           # position y
    color = Tier,       # colour
    size = Amount,      # size
    shape = CarType)) + # shape
  geom_point(alpha = 0.4) +
  scale_color_viridis_c() +
  scale_size(range = c(1, 8)) +
  scale_shape_manual(values =
    c(N = 16, U = 17)) +
  theme_minimal() +
  labs(title = "5 Variables,
    1 Scatter")

# Marginal distributions
library(ggExtra)
ggMarginal(p, type = "histogram")
\end{lstlisting}
    \end{column}
    \begin{column}{0.48\textwidth}
      \textbf{Python}
\begin{lstlisting}[style=pythonstyle]
fig, ax = plt.subplots()
for ct, marker in [("N","o"),("U","s")]:
    mask = ecar["CarType"] == ct
    sc = ax.scatter(
        ecar.loc[mask, "FICO"],
        ecar.loc[mask, "Rate"],
        c=ecar.loc[mask, "Tier"],
        s=ecar.loc[mask, "Amount"]/300,
        marker=marker,
        alpha=0.4, cmap="viridis",
        label=ct)
plt.colorbar(sc, label="Tier")
ax.legend(title="Car Type")

# Marginal: seaborn jointplot
sns.jointplot(data=ecar,
    x="FICO", y="Rate",
    hue="CarType",
    kind="scatter", height=6)
\end{lstlisting}
    \end{column}
  \end{columns}
\end{frame}

\section{Aesthetic Overload}

\begin{frame}{More Encodings $\neq$ More Information}
  \begin{center}
    \includegraphics[width=0.92\textwidth]{figures_w05m01/overload}
  \end{center}
  \begin{alertblock}{Pro-Tip}
    The human visual system can track $\sim$4 encoding channels
    simultaneously (Healey \& Enns, 2012). Beyond that, the chart
    becomes a puzzle rather than a communication tool.
    \textbf{Cognitive limit}: $\leq$8 colour hues, $\leq$6 shapes,
    $\leq$5 size levels. If you need more groups, \textbf{facet}
    instead of encoding.
  \end{alertblock}
\end{frame}

\section{Five Strategies}

\begin{frame}{Strategies for High-Dimensional Visualization}
  \begin{center}
    \includegraphics[width=0.95\textwidth]{figures_w05m01/strategies}
  \end{center}
\end{frame}

\begin{frame}{Choosing a Strategy}
  \centering
  \begin{tabular}{lll}
    \toprule
    \textbf{Situation} & \textbf{Strategy} & \textbf{Module} \\
    \midrule
    $\leq$5 variables & Multi-encoded scatter & This module \\
    5--10 variables & Parallel coordinates & M02 \\
    5--10 variables (ranks) & Radar / spider chart & M02 \\
    10--100 variables & PCA / t-SNE / UMAP & M03 \\
    Variables $\times$ observations matrix & Heatmap + clustering & M04 \\
    Geographic data & Maps (choropleth, point) & M05--M08 \\
    $>$5 categorical groups & Small multiples (facet) & W03-M08 \\
    \bottomrule
  \end{tabular}
\end{frame}

\section{Aggregation for Large $n$}

\begin{frame}{When $n$ Is Too Large for Scatter}
  \begin{center}
    \includegraphics[width=0.88\textwidth]{figures_w05m01/aggregation}
  \end{center}
\end{frame}

\begin{frame}[fragile]{Aggregation in Practice}
  \begin{columns}[T]
    \begin{column}{0.48\textwidth}
      \textbf{R}
\begin{lstlisting}[style=rstyle]
# Hexbin (for 50K+ points)
ggplot(ecar, aes(x = FICO,
    y = Rate)) +
  geom_hex(bins = 50) +
  scale_fill_viridis_c() +
  theme_minimal()

# Sample + scatter
ecar |> slice_sample(n = 2000) |>
  ggplot(aes(x = FICO, y = Rate)) +
  geom_point(alpha = 0.1, size=0.5)

# Smooth (LOESS through cloud)
ggplot(ecar, aes(x=FICO, y=Rate)) +
  geom_smooth() +
  theme_minimal()

# Marginal histograms
library(ggExtra)
p <- ggplot(ecar, aes(x=FICO,
    y=Rate)) + geom_point(alpha=0.05)
ggMarginal(p, type = "histogram")
\end{lstlisting}
    \end{column}
    \begin{column}{0.48\textwidth}
      \textbf{Python}
\begin{lstlisting}[style=pythonstyle]
# Hexbin
ax.hexbin(ecar["FICO"], ecar["Rate"],
    gridsize=50, cmap="viridis",
    mincnt=1)
plt.colorbar(label="Count")

# Sample + scatter
sub = ecar.sample(2000)
ax.scatter(sub["FICO"], sub["Rate"],
    s=3, alpha=0.1)

# 2D KDE
sns.kdeplot(data=ecar, x="FICO",
    y="Rate", fill=True,
    cmap="Blues", levels=10)

# Jointplot (marginals included)
sns.jointplot(data=ecar,
    x="FICO", y="Rate",
    kind="hex", height=6)
\end{lstlisting}
    \end{column}
  \end{columns}
\end{frame}

\section{Synthesis}

\begin{frame}{Key Takeaways}
  \begin{enumerate}
    \item The \textbf{curse of dimensionality}: same $n$ in more
          dimensions = sparser coverage. Standard scatter fails
          beyond $\sim$3 variables.
    \item \textbf{Encoding channels} are ranked: position $>$ length
          $>$ angle $>$ area $>$ colour hue $>$ saturation $>$ shape.
          Use the highest-ranked channels for the most important variables.
    \item \textbf{Budget}: a single scatter can encode $\sim$5 variables
          (x, y, colour, size, shape). Beyond that, \textbf{facet}
          or switch chart types.
    \item \textbf{Aesthetic overload} ($>$8 colours, $>$6 shapes)
          overwhelms the visual system. Simpler $>$ denser.
    \item Five strategies: \textbf{reduce} (PCA), \textbf{encode} (multi-scatter),
          \textbf{facet} (small multiples), \textbf{parallel} coordinates,
          \textbf{heatmap} + clustering.
    \item For large $n$: \textbf{hexbin}, \textbf{sample}, \textbf{smooth},
          \textbf{2D KDE}, \textbf{marginal} distributions.
  \end{enumerate}
\end{frame}

\begin{frame}{Essential Readings}
  \begin{itemize}
    \item Cleveland, W.~S. \& McGill, R. (1984). ``Graphical Perception.''
          \emph{JASA}, 79(387). (Encoding channel ranking.)
    \item Healey, C.~G. \& Enns, J.~T. (2012). ``Attention and Visual
          Memory in Visualization.'' \emph{IEEE TVCG}, 18(7).
    \item Munzner, T. (2014). \emph{Visualization Analysis and Design}.
          CRC Press. (Channel effectiveness, task taxonomy.)
    \item Wilke, C.~O. (2019). \emph{Fundamentals of Data Visualization},
          Chapter 2 (Encoding Channels).
  \end{itemize}
  \vspace{6pt}
  \textbf{Next Module}: Parallel Coordinates \& Radar Charts (W05-M02)
\end{frame}

\end{document}
